% titel: Blatt 9, A3: Dampfdruck einer idealen Lösung (Raoult)
% skript: Mortimer Kap. 13 (Raoultsches Gesetz)

% Teil I: Werkzeuge, kein Lösungsweg.

Bei einer idealen Lösung stellt sich der Gesamtdampfdruck additiv aus den
Dampfdrücken der reinen Komponenten zusammen.

\begin{werkzeugkasten}[Raoultsches Gesetz]
Jede Komponente $i$ trägt zum Gesamtdampfdruck bei:
\[
  p_i = x_i \cdot p_i^\circ,
\]
mit Molenbruch $x_i$ und Dampfdruck $p_i^\circ$ der \emph{reinen}
Komponente. Der Gesamtdruck über einer Mischung aus zwei flüchtigen Stoffen
A und B:
\[
  p_{\text{ges}} = x_A p_A^\circ + x_B p_B^\circ,\qquad
  x_A + x_B = 1 .
\]
\textit{Skript:} Mortimer Kap. 13 (ideale Lösungen, Raoult).
\end{werkzeugkasten}

\begin{achtung}[Darauf achten]
\begin{itemize}
  \item Sieden: die Lösung siedet, sobald $p_{\text{ges}}$ den Aussendruck
        (hier 1 bar) erreicht.
  \item $x_A+x_B=1$ verwenden, um auf eine Unbekannte zu reduzieren.
  \item Nach den \emph{Stoffmengenanteilen} (Molenbrüchen) ist gefragt,
        nicht nach Massen.
\end{itemize}
\end{achtung}
